\documentclass[a4paper,11pt]{book}
%\documentclass[a4paper,twoside,11pt,titlepage]{book}
\usepackage{listings}
\usepackage[utf8]{inputenc}
\usepackage[spanish]{babel}

% \usepackage[style=list, number=none]{glossary} %
%\usepackage{titlesec}
%\usepackage{pailatino}
\usepackage[table,xcdraw]{xcolor}
%\decimalpoint
\usepackage{dcolumn}
\usepackage{float}
\newcolumntype{.}{D{.}{\esperiod}{-1}}
\makeatletter
%\addto\shorthandsspanish{\let\esperiod\es@period@code}
\makeatother


%\usepackage[chapter]{algorithm}
\RequirePackage{verbatim}
%\RequirePackage[Glenn]{fncychap}
\usepackage{fancyhdr}
\usepackage{graphicx}
\usepackage{afterpage}

\usepackage{longtable}

\usepackage[pdfborder={000}]{hyperref} %referencia

% ********************************************************************
% Re-usable information
% ********************************************************************
\newcommand{\myTitle}{Trabajo Fin de Máster\xspace}
\newcommand{\myDegree}{Máster Universitario de Investigación en Ingeniería de Software y Sistemas Informáticos\xspace}
\newcommand{\myName}{César Hugo Bárzano Cruz\xspace}
\newcommand{\myProf}{Rubén Heradio\xspace}
\newcommand{\myFaculty}{ Universidad Nacional de Educación a Distancia\xspace}
\newcommand{\myFacultyShort}{UNED-Facultad de informática\xspace}
\newcommand{\myDepartment}{\xspace}
\newcommand{\myUni}{\protect{ Universidad Nacional de Educación a Distancia}\xspace}
\newcommand{\myLocation}{Madrid\xspace}
\newcommand{\myTime}{\today\xspace}
\newcommand{\myVersion}{Version 0.1\xspace}


\hypersetup{
pdfauthor = {\myName hugobarzano@gmail.com},
pdftitle = {\myTitle},
pdfsubject = {},
pdfkeywords = {},
pdfcreator = {LaTeX con el paquete TEXmaker},
pdfproducer = {pdflatex}
}

%\hyphenation{}


%\usepackage{doxygen/doxygen}
%\usepackage{pdfpages}
\usepackage{url}
\usepackage{colortbl,longtable}
\usepackage[stable]{footmisc}
%\usepackage{index}

%\makeindex
%\usepackage[style=long, cols=2,border=plain,toc=true,number=none]{glossary}
% \makeglossary

% Definición de comandos que me son tiles:
%\renewcommand{\indexname}{Índice alfabético}
%\renewcommand{\glossaryname}{Glosario}

\pagestyle{fancy}
\fancyhf{}
\fancyhead[LO]{\leftmark}
\fancyhead[RE]{\rightmark}
\fancyhead[RO,LE]{\textbf{\thepage}}
\renewcommand{\chaptermark}[1]{\markboth{\textbf{#1}}{}}
\renewcommand{\sectionmark}[1]{\markright{\textbf{\thesection. #1}}}

\setlength{\headheight}{1.5\headheight}


\usepackage{bera}% optional: just to have a nice mono-spaced font
\usepackage{listings}
\usepackage{xcolor}


\colorlet{punct}{red!60!black}
\definecolor{background}{HTML}{EEEEEE}
\definecolor{delim}{RGB}{20,105,176}
\colorlet{numb}{magenta!60!black}

\lstdefinelanguage{json}{
    basicstyle=\normalfont\ttfamily,
    numbers=left,
    numberstyle=\scriptsize,
    stepnumber=1,
    numbersep=8pt,
    showstringspaces=false,
    breaklines=true,
    frame=lines,
    backgroundcolor=\color{background},
    literate=
     *{0}{{{\color{numb}0}}}{1}
      {1}{{{\color{numb}1}}}{1}
      {2}{{{\color{numb}2}}}{1}
      {3}{{{\color{numb}3}}}{1}
      {4}{{{\color{numb}4}}}{1}
      {5}{{{\color{numb}5}}}{1}
      {6}{{{\color{numb}6}}}{1}
      {7}{{{\color{numb}7}}}{1}
      {8}{{{\color{numb}8}}}{1}
      {9}{{{\color{numb}9}}}{1}
      {:}{{{\color{punct}{:}}}}{1}
      {,}{{{\color{punct}{,}}}}{1}
      {\{}{{{\color{delim}{\{}}}}{1}
      {\}}{{{\color{delim}{\}}}}}{1}
      {[}{{{\color{delim}{[}}}}{1}
      {]}{{{\color{delim}{]}}}}{1},
}


\newcommand{\HRule}{\rule{\linewidth}{0.5mm}}
%Definimos los tipos teorema, ejemplo y definición podremos usar estos tipos
%simplemente poniendo \begin{teorema} \end{teorema} ...
\newtheorem{teorema}{Teorema}[chapter]
\newtheorem{ejemplo}{Ejemplo}[chapter]
\newtheorem{definicion}{Definición}[chapter]

\definecolor{gray97}{gray}{.97}
\definecolor{gray75}{gray}{.75}
\definecolor{gray45}{gray}{.45}
\definecolor{gray30}{gray}{.94}

\lstset{ frame=Ltb,
     framerule=0.5pt,
     aboveskip=0.5cm,
     framextopmargin=3pt,
     framexbottommargin=3pt,
     framexleftmargin=0.1cm,
     framesep=0pt,
     rulesep=.4pt,
     backgroundcolor=\color{gray97},
     rulesepcolor=\color{black},
     %
     stringstyle=\ttfamily,
     showstringspaces = false,
     basicstyle=\scriptsize\ttfamily,
     commentstyle=\color{gray45},
     keywordstyle=\bfseries,
     %
     numbers=left,
     numbersep=6pt,
     numberstyle=\tiny,
     numberfirstline = false,
     breaklines=true,
   }

% minimizar fragmentado de listados
\lstnewenvironment{listing}[1][]
   {\lstset{#1}\pagebreak[0]}{\pagebreak[0]}

\lstdefinestyle{CodigoC}
   {
	basicstyle=\scriptsize,
	frame=single,
	language=C,
	numbers=left
   }
\lstdefinestyle{CodigoC++}
   {
	basicstyle=\small,
	frame=single,
	backgroundcolor=\color{gray30},
	language=C++,
	numbers=left
   }


\lstdefinestyle{Consola}
   {basicstyle=\scriptsize\bf\ttfamily,
    backgroundcolor=\color{gray30},
    frame=single,
    numbers=none
   }


\newcommand{\bigrule}{\titlerule[0.5mm]}
\usepackage{enumitem}


%Para conseguir que en las páginas en blanco no ponga cabecerass
\makeatletter
\def\clearpage{%
  \ifvmode
    \ifnum \@dbltopnum =\m@ne
      \ifdim \pagetotal <\topskip
        \hbox{}
      \fi
    \fi
  \fi
  \newpage
  \thispagestyle{empty}
  \write\m@ne{}
  \vbox{}
  \penalty -\@Mi
}
\makeatother

\usepackage{pdfpages}
\begin{document}
\begin{titlepage}
 
 
\newlength{\centeroffset}
\setlength{\centeroffset}{-0.5\oddsidemargin}
\addtolength{\centeroffset}{0.5\evensidemargin}
\thispagestyle{empty}

\noindent\hspace*{\centeroffset}\begin{minipage}{\textwidth}

\centering
\includegraphics[width=0.7\textwidth]{imagenes/Logo-uned.jpg}\\[1.1cm]


{\Huge\bfseries Máster Universitario de Investigación en Ingeniería De Software Y Sistemas Informáticos\\
}
\noindent\rule[-1ex]{\textwidth}{3pt}\\[3.5ex]
{\large\bfseries 31105151 - Trabajo Fin de Máster}
\end{minipage}

\vspace{2.5cm}
\noindent\hspace*{\centeroffset}\begin{minipage}{\textwidth}
\centering

\textbf{Autor}\\ {César Hugo Bárzano Cruz}\\[2.5ex]
\textbf{Tutor}\\ {Rubén Heradio}\\[2.5ex]


%\includegraphics[width=0.3\textwidth]{imagenes/Logo-master.png}\\[0.1cm]
\textsc{Generative Cloud Manager: code-runner}\\
\textsc{---}\\
2019/2020
\end{minipage}
%\addtolength{\textwidth}{\centeroffset}
%\vspace{\stretch{2}}
\end{titlepage}

\begin{titlepage}
 
 

\setlength{\centeroffset}{-0.5\oddsidemargin}
\addtolength{\centeroffset}{0.5\evensidemargin}
\thispagestyle{empty}

\noindent\hspace*{\centeroffset}\begin{minipage}{\textwidth}

\centering


{\Huge\bfseries Máster Universitario de Investigación en Ingeniería De Software Y Sistemas Informáticos\\}

\noindent\rule[-1ex]{\textwidth}{3pt}\\[3.5ex]
{\large\bfseries 31105151 - Trabajo Fin de Máster}
\textbf{Generative Cloud Manager: code-runner}\\[2.5ex]
{Consultar el tipo de trabajo: Trabajo Propuesto por Alumno}\\[2.5ex]
\end{minipage}

\vspace{2.5cm}
\noindent\hspace*{\centeroffset}\begin{minipage}{\textwidth}
\centering

\textbf{Autor}\\ {César Hugo Bárzano Cruz}\\[2.5ex]
\textbf{Tutor}\\ {Rubén Heradio}\\[2.5ex]


\includegraphics[width=0.3\textwidth]{imagenes/Logo-master.png}\\[0.1cm]
\textsc{---}\\
2019/2020
\end{minipage}

\end{titlepage}




\vspace{0.7cm}
\noindent{\textbf{Resumen}}\\


Este proyecto demuestra como el uso de conceptos ligados a la  Programación Generativa y la integración de diversos servicios web pueden utilizarse de manera conjunta para sintetizar, generalizar y automatizar el proceso de desarrollo y operación de aplicaciones software. La lectura aquí propuesta, junto con el sistema software resultante abordan  el proceso de desarrollo de aplicaciones software mediante un enfoque generativo, con esto  se pretende proporcionar al desarrollador software de un sistema web orientado a la generación, construcción y despliegue de aplicaciones software, concretamente de aplicaciones  web o servicios de datos en la nube.
 
 En este proyecto se considera que las etapas que forman el ciclo de vida o proceso de desarrollo de aplicaciones software está altamente segmentado, por lo que presenta una solución generativa y automatizada al proceso de desarrollo software. Con este proyecto no se pretende desprestigiar a los procesos de desarrollo software tradicionales, lo que intenta es demostrar como la unificación de los conceptos comunes entre aplicaciones web de diversa naturaleza así como la síntesis y  automatización genérica de las etapas básicas de producción software pueden dar lugar a soluciones funcionales ahorrando al desarrollador numerosas horas de trabajo. De igual forma, el sistema resultante  puede ser aprovechado por el desarrollador inexperto como punto de partida en su aprendizaje de las denominadas tecnologías de la información.
 
 
Este proyecto pone de manifiesto que ante la gran diversidad de tecnologías y soluciones ofrecidas por la Cloud Computing  existen definiciones, conceptos y  etapas básicas que son  comunes  y transversales  a todas ellas, independientemente de su naturaleza, dominio y propósito. Este proyecto sintetiza el proceso de desarrollo software en 3 etapas fundamentales: 

\begin{enumerate}
\item \textbf{ Code Generation. }  Etapa inicial donde se sintetiza el proceso de desarrollo software, es decir, se produce el código necesario para la solución especificada por el desarrollador. 
\item \textbf{ Artefact Generation, }  El objetivo de esta etapa es aprovechar el código fuente generado en la etapa Code Generation para producir un artefacto auto-contenido junto con la configuración necesaria para su ejecución en un entorno virtual. 
\item \textbf{ Deploy Generation. } En esta etapa se aprovecha el producto de la etapa Artefact Generation para disponibilizar la solución en Internet. 
\end{enumerate}
 

En los siguientes capítulos se profundiza en estas ideas, detallando cuales son los objetivos de alto nivel o casos de uso que se esperan del sistema. Se describirá el problema planteado y se analizará en detalle la solución alcanzada.\\ \\

\noindent{\textbf{Palabras Clave}: Cloud Computing, Programación Generativa, Desarrollo Software, Código Fuente, Automatización, Integración, Despliegue, Aplicaciones, Servicios, Contenedores.}\\

\afterpage{\null\newpage}
\newpage



%\frontmatter
\tableofcontents
\listoffigures
%\listoftables

%
%\mainmatter
%\setlength{\parskip}{5pt}

%\input{capitulos/01_Introduccion}



\chapter{Introducción}

\section{Motivación}

La motivación de este proyecto es la de mejorar el proceso de desarrollo de aplicaciones web y servicios cloud, ayudando al desarrollador a comprender y adquirir conocimientos relativos a las tecnologías de la información.  En el proceso de desarrollo de software informático no existe un proceso universal que indique las pautas a seguir para la correcta producción software. Comúnmente, este proceso ha de amoldarse a la naturaleza del software que se desea desarrollar, a la tecnología con la que ha de desarrollarse y como no a los recursos humanos o desarrolladores que van a construir dicho software. La motivación de este proyecto es la de dotar a los desarroladores de un sistema que los abstraiga de estas etapas, premiando la rápida producción y disponibilidad del software deseado. 

\section{Objetivos}

El objetivo de este proyecto es el de sintetizar el proceso de desarrollo y operación de software en un conjunto automatizado de etapas. Dicha síntesis consiste en generalizar la naturaleza de los distintos sistemas de información que pueden ser consumidos a través de Internet en una única entidad denominada comúnmente como aplicación. La generalización ocurre al manejar los distintos sistemas de información:

\begin{enumerate}
\item \textbf{ Aplicaciones web estáticas }  
\item \textbf{ Aplicaciones web dinámicas}
\item \textbf{ Interfaces web }  
\item \textbf{ Api/Rest }  
\item \textbf{ Servicios }  
\item \textbf{ Cualquier pieza software que pueda ser consumida mediante protocolo tcp/http }  
\end{enumerate}

como una única  entidad \textbf{aplicación} y  someter a esta entidad al mismo proceso o etapas independientemente de cual sea su naturaleza o dominio cumpliendo con el objetivo de dada una especificación de aplicación, producir el código y/o configuración necesaria para materializar dicha especificación en un producto software, construir un artefacto auto-contenido que permita la ejecución del software generado, ejecutar dicho artefacto sobre la infraestructura virtual necesaria para que la aplicación especificada sea accesible por Internet y que de esta forma el usuario de este proyecto sea capaz de consumir la aplicación generada.  
 
El objetivo de este proyecto es el de dotar a los usuarios o desarrolladores de un  sistema  generativo y automatizado de aplicaciones que mejore el proceso de desarrollo software y facilite el aprendizaje de nuevas tecnologías. El principal objetivo del proyecto es acuñar el lema "From Spec to Cloud" es decir, dada una especificación de aplicación de alto nivel, materializarla en un sistema de información disponible en la nube. 

 

\textbf{OBJ-1.} 

\textbf{OBJ-2.}

\chapter{Descripción del Problema}

\section{El Problema}

El auge de las tecnologías de la información a nivel empresarial, social, etico...blacbla


La cloud esta formada por un conjunto de tecnologías que habilitan al consumidor una sería de servicios de información para ser consumidos vía Internet. Cada vez son mas y mas las empresas y entidades confían en este tipo de servicios para realizar sus actividades económicas. La diversidad de soluciones demandadas ha creado un ecosistema en el que interactuan diversos sistemas y diversas tecnologías. Debido a esto la demanda de personal cualificado o desarrolladores también se ha incrementado. 

Hoy en día la sociedad esta rodeada de 


Diversidad de aplicaciones cloud. Nuevas tecnologías entorno variado y cambiante. Demanda de nuevas tecnologías.
Capacitación para desarrolladores. Diversidad de aplicaciones, mismo tratamiento, simplificar proceso de aprendizaje y desarrollo software. FullStack. 


\section{Generadores de aplicaciones}
\section{Despliegue de código}
 

\chapter{Descripción de la Solución}


Con el objetivo de generalizar y definir aplicación software se han decidido usar las partes:

Repositorio
Readme
Makefile
Licencia
src codigo fuente
config configuración


\begin{enumerate}
\item \textbf{ Code Generation }  Etapa inicial donde se sintetiza el proceso de desarrollo software, es decir se produce el código necesario para la solución especificada por el desarrollador. 
\item \textbf{ Artefact Generation }  El objetivo de esta etapa es aprovechar el código fuente generado en la etapa Code Generation para producir un artefacto auto-contenido junto con la configuración necesaria para su ejecución en un entorno virtual. 
\item \textbf{ Deploy Generation } En esta etapa se aprovecha el producto de la etapa Artefact Generation para disponibilizar la solución en Internet. 
\end{enumerate}


Programación generativa: generar, automatizar  y despliegue accesible a Internet.


La aplicación software es el elemento de negocio de este sistema.  Para materializar esto, el sistema crea un repositorio de código por cada aplicación creada para albergar la producción de código especificado por  el desarrolador. 


Generar: Aplicaciones completas, independientemente de la tecnología, utilidad o fin. Las aplicaciones generadas por el sistema
son auto-contenidas y automatizadas, es decir mantienen una semántica común 




d
test
run
pull
push

Lo que permite al desarrollador integrar fácilmente nuevos cambios. 

Definir eel permino produción dentro del marco de este sistema, consiste en la siguiente serie de etapas:



\section{Análisis}

Como se ha visto, el termino producción en el termino de este sistema considera 

\subsection{Requisitos de Información }
Los requisitos de información se caracterizan por reunir la información relevante para la solución, que debe gestionar y almacenar el sistema software.\\

\textbf{RI-1. Workspace:} Representación del espacio de trabajo del usuario. Contiene el conjunto de las aplicaciones que el usuario ha creado usando el sistema. Concepto similar al del Escritorio o "Desktop" de un sistema operativo. 
Contenido: nombre de usuario o propietario del workspace, referencia al conjunto de aplicaciones generadas.


\textbf{RI-2. App:} Representación de cada una de las aplicaciones generadas por el sistema. 
Contenido:
\begin{lstlisting}[language=json,firstnumber=1]
{
    "_id" : "appName",
    "repo" : "https://github.com/user/appName.git",
    "spec" : {
        "dockerId" : "appNameDockerId",
        "port" : "4343",
        "nature" : "staticApp"
    },
    "des" : "",
    "url" : "http://deploy.domain:4343",
    "owner" : "user>s",
    "status" : "running"
}
\end{lstlisting}

\textbf{RI-3. Información de Usuario:} Información sensible del usuario o desarrollador que va a usar el sistema. Credenciales al sistema de gestión de repositorios para el código fuente, en este caso relativa al perfil de github. 
Contenido: nombre de usuario, correo electrónico, organización. \\


\subsection{Requisitos Funcionales }
Como se define en la ingeniería de requisitos, los requisitos funcionales establecen el comportamientos del sistema.\\

\textbf{RF-1. Generación de aplicaciones:} El sistema ha de generar  aplicaciones dada una especificación.\\
   

	\textbf{RF-1.1.} El sistema generará aplicaciones de diversa naturaleza, tecnología o dominio 
	\textbf{RF-1.2.}  El sistema creará el repositorio de código necesario para albergar las aplicaciones generadas
	RF-1.2. El sistema permitirá eliminar aplicaciones.

	RF-1.3. El sistema permitirá ejecutar aplicaciones.

	RF-1.4. El sistema permitirá detener aplicaciones.

	RF-1.5. El sistema permitirá visualizar información de las aplicaciones.\\



\textbf{RF-1. Gestión de aplicaciones:} El sistema ha de gestionar el ciclo de vida de las aplicaciones generadas.\\
   

	\textbf{RF-1.1.} El sistema permitirá crear aplicaciones.
	
	\quad \textbf{RF-1.1.1.}El sistema generará aplicaciones de diversa naturaleza, tecnología o dominio.
	
	\quad \textbf{RF-1.1.2.}El sistema creará y  guardará en repositorios de código las aplicaciones generadas.
	
	\quad \textbf{RF-1.1.3.}El sistema construirá los artefactos necesarios para ejecutar las aplicaciones generadas.
	
	\quad \textbf{RF-1.1.4.} El sistema ejecutará los artefactos construidos para disponibilizar vía Internet las aplicaciones creadas. \\



	\textbf{RF-1.2.} El sistema permitirá eliminar aplicaciones.
	
	\quad \textbf{RF-1.2.1.}El sistema detendrá los artefactos en ejecución de aplicaciones eliminadas. 
	
	\quad \textbf{RF-1.2.2.}El sistema eliminará los artefactos construidos de aplicaciones eliminadas. 
	
	\quad \textbf{RF-1.2.3.}El sistema eliminará los repositorios de código de aplicaciones eliminadas. \\
	

	
	\textbf{RF-1.3.} El sistema permitirá ejecutar aplicaciones.
	
	\quad \textbf{RF-1.3.1.} El sistema será capaz de hacer a las aplicaciones accesibles vía Internet.  
	
	\quad \textbf{RF-1.3.2.} El sistema proporcionará las herramientas  y configuraciones necesarias para que las aplicaciones puedan ser ejecutas en local y ser accesibles vía l\textbf{ocalhost}  \\
	

	\textbf{RF-1.4.} El sistema permitirá detener aplicaciones.
	
		\quad \textbf{RF-1.4.1.} El sistema será capaz de detener las aplicaciones que sean accesibles vía Internet. 
		
			\quad \textbf{RF-1.4.2.} El sistema proporcionará las herramientas  y configuraciones necesarias para que las aplicaciones puedan ejecutas en local sean detenidas. \\
			
	
	\textbf{RF-1.5.} El sistema permitirá visualizar información de las aplicaciones.\\
	
		\quad \textbf{RF-1.5.1.} El sistema permitirá consultar  información de utilidad relativa a las aplicaciones,  como por ejemplo, especificación, URL del repositorio donde se alberga el código fuente, URL oo IP donde la aplicación puede ser consumida vía Internet, etc...
	
			\quad \textbf{RF-1.5.2.} El sistema permitirá consultar información de DEBUG de las aplicaciones accesibles vía Internet, es decir, permitirá consultar LOGs de aplicaciones en ejecución. 
			
	\textbf{RF-1.6.} El sistema permitirá visualizar información del estado del ciclo de vida de las aplicaciones. El ciclo de vida de las aplicaciones del sistema se representará con el siguiente flujo de estados:\\
			
\textbf{RF-2. Gestión de credenciales:} El sistema proporcionará las credenciales necesarias para trabajar con las aplicaciones .\\
   

	RF-2.1. El sistema permitirá 


\subsection{Requisitos No Funcionales }
Los requisitos no funcionales, se refieren a todos los requisitos que no describen información a guardar, ni funciones a realizar por el sistema, sino características de funcionamiento.\\


\textbf{RNF-1} Se necesitará acceso a Internet para utilizar las funcionalidades del sistema.\\
\textbf{RNF-2} Se necesitará una cuenta en github para utilizar las funcionalidades del sistema.\\


\subsection{Casos de Uso}
Los diagramas de casos de uso, son diagramas UML que representan gráficamente a todos los elementos que forman parte del modelo de casos de uso junto con la frontera del sistema. Delimitan el sistema a diseñar. Determinan el contexto del uso del sistema. Describen el punto de vista de los usuarios  en el sistema.



\chapter{ Medición y Evaluación}

\chapter{Conclusiones}




\begin{thebibliography}{aaaa}



\bibitem[1]{pt1} \textsc{Tema 1},
\textit{ pag 8. Visión por computador: imágenes digitales y aplicaciones. Gonzalo Pajares Martinsanz y Jesús Manuel de la Cruz García}




\end{thebibliography}

\chapter{ Anexo}

\section{Plantilla Corrección Tipo Test }




%
%
%%\nocite{*}
%\bibliography{bibliografia/bibliografia}\addcontentsline{toc}{chapter}{Bibliografía}
%\bibliographystyle{miunsrturl}
%
%\appendix

%\input{apendices/manual_usuario/manual_usuario}
%%\input{apendices/paper/paper}
%\input{glosario/entradas_glosario}
% \addcontentsline{toc}{chapter}{Glosario}
% \printglossary

\thispagestyle{empty}

\end{document}
